\section{IoT Internetworking}

\subsection{Motivation}

The recorded lecture is provided \myhref{https://www.youtube.com/watch?v=9YRa5QyNf4w&feature=youtu.be}{here}.

\begin{ejercicio}
    Why do we need new protocols for the IoT instead of using established Internet protocols? Describe one challenge existing protocols face.\\

    New protocols are needed for the IoT because existing Internet protocols, such as TCP/IP, are not optimized for the resource-constrained nature of IoT devices. One challenge existing protocols face is that they may require larger buffer sizes than what is available on IoT devices, which can lead to performance issues and increased latency.
\end{ejercicio}

\begin{ejercicio}
    Which central protocol can we not replace (completely) for the IoT? Why is it harder to replace it than others?\\

    The central protocol that we cannot replace (completely) for the IoT is the Internet Protocol (IP). It is harder to replace IP than other protocols because it is the fundamental protocol that enables communication and data exchange between different devices and applications on the Internet. It is the backbone of the Internet and, technically speaking, being connected to the Internet requires having an IP address and using IP for communication. Therefore, while we can use different protocols for the application layer or the transport layer, we still need to use IP at the network layer in order to enable connectivity and communication with the Internet.
\end{ejercicio}

\begin{ejercicio}
    We introduced two approaches for Smart Things to connect to the Internet. Name these approaches, describe them briefly and give one advantage and one disadvantage of each, compared to the other.\\

    The two approaches for Smart Things to connect to the Internet are:
    \begin{itemize}
        \item \ul{Application layer proxy}: In this approach, there is a proxy (usually implemented in the IoT gateway) that translates between the protocols used by the IoT Things and the protocols used by the Internet. The IoT Thing communicates with the proxy using dedicated IoT protocols that have nothing to do with the Internet protocols, and the proxy stores the data received from the IoT Things and makes it available to the Internet using standard Internet protocols.
        \begin{itemize}
            \item Advantage: It allows for the use of dedicated IoT protocols that are optimized for the resource-constrained nature of IoT devices, which can lead to better performance and energy efficiency.
            \item Disadvantage: It introduces an additional layer of complexity and can lead to increased latency and reduced scalability (due to the dedicated software), as all communication between the IoT Things and the Internet needs to go through the proxy.
        \end{itemize}

        \item \ul{IP Router}: In this approach, the IoT Things directly use IP, but with some optimizations and adaptations to address the challenges and limitations mentioned above. There is no need for specific apps or software to interact with the IoT Things, as they can be accessed using standard Internet protocols and interfaces.
        \begin{itemize}
            \item Advantage: It provides better scalability and interoperability, as the IoT Things can communicate directly with other devices and applications on the Internet without the need for a proxy.
            \item Disadvantage: It may require more complex and resource-intensive protocols and technologies to enable efficient communication and data exchange between the IoT Things and the Internet, which can lead to increased power consumption and reduced performance for resource-constrained IoT devices.
        \end{itemize}
    \end{itemize}
\end{ejercicio}

\subsection{IEEE 802.15.4}

The recorded lectures are provided \myhref{https://www.youtube.com/watch?v=utk_hszYO_o}{here (Part A)} and \myhref{https://www.youtube.com/watch?v=dOYCUwEpEnw}{here (Part B)}.

\begin{ejercicio}
    What are the differences between the PHY and the MAC layer?
\end{ejercicio}

\begin{ejercicio}
    Discuss pros and cons of using 16 bit addresses for 802.15.4 instead of 64 bit addresses.
\end{ejercicio}

\begin{ejercicio}
    Describe the basic idea of CSMA/CA and compare it to frequency hopping. You may use a diagram / drawing to describe your answer.
\end{ejercicio}

\begin{ejercicio}
    Briefly describe and compare TDMA and FDMA. Which one is better?
\end{ejercicio}

\begin{ejercicio}
    An 802.15.4 network is deployed in the following three application scenarios. Which CSMA/CA MAC Layer version would you use? Justify your answer (e.g., ``In this case, we should use slotted CSMA/CA, because...'').
    \begin{itemize}
        \item A smart home network, where only 10 temperature sensors are installed. Each sensor reports a measurement once per minute.
        \item A smart clinic network that transmits monitoring data about the patients and alarm notifications.
    \end{itemize}
\end{ejercicio}

\begin{ejercicio}
    When unslotted CSMA/CA is applied, at each trail the back-off window size will be doubled. However, it doesn't always lead to a longer delay, why?
\end{ejercicio}